\documentclass[11pt,a4paper]{article}

\usepackage[utf8]{inputenc}
\usepackage[T1]{fontenc}
\usepackage{lmodern}
\usepackage[margin=2.5cm]{geometry}
\usepackage{hyperref}
\usepackage{xcolor}
\usepackage{listings}
\usepackage{booktabs}
\usepackage{enumitem}
\usepackage{parskip}
\usepackage{amssymb}
\usepackage{fancyhdr}
\usepackage{titlesec}

% ── Colours ──────────────────────────────────────────────────────────────────
\definecolor{cbaKeyword}{RGB}{0,100,170}
\definecolor{cbaString}{RGB}{160,32,32}
\definecolor{cbaComment}{RGB}{100,100,100}
\definecolor{cbaOp}{RGB}{140,70,20}
\definecolor{cbaBg}{RGB}{248,248,248}
\definecolor{cbaFrame}{RGB}{200,200,200}
\definecolor{linkcolor}{RGB}{0,80,160}

% ── Hyperref setup ───────────────────────────────────────────────────────────
\hypersetup{
  colorlinks=true,
  linkcolor=linkcolor,
  urlcolor=linkcolor,
  pdftitle={Cambria Language Reference},
  pdfauthor={},
}

% ── Listings: Cambria language ───────────────────────────────────────────────
\lstdefinelanguage{cambria}{
  morekeywords={fun,rec,handler,return,finally,do,in,if,then,else,with,handle,
                inl,inr,case,of,effect,Unit,Void,Int,Bool,Double,Str,Name,Map,
                List,true,false},
  sensitive=true,
  morecomment=[l]{--},
  morestring=[b]",
  literate=
    {->}{{$\to$}}2
    {<-}{{$\leftarrow$}}2
    {~>}{{$\leadsto$}}2
    {=>}{{$\Rightarrow$}}2
    {++}{{+\!+}}2
    {::}{{:\!:}}2
    {/=}{{$\neq$}}1
    {==}{{$==$}}2
    {<=}{{$\leq$}}2
    {>=}{{$\geq$}}2
    {\\}{{$\lambda$}}1
    {()}{{()}}2
    {[]}{{[\,]}}2,
}

\lstdefinestyle{cba}{
  language=cambria,
  basicstyle=\ttfamily\small,
  keywordstyle=\color{cbaKeyword}\bfseries,
  stringstyle=\color{cbaString},
  commentstyle=\color{cbaComment}\itshape,
  backgroundcolor=\color{cbaBg},
  frame=single,
  rulecolor=\color{cbaFrame},
  framesep=6pt,
  xleftmargin=8pt,
  xrightmargin=8pt,
  breaklines=true,
  breakatwhitespace=true,
  showstringspaces=false,
  tabsize=2,
  columns=flexible,
  keepspaces=true,
  aboveskip=10pt,
  belowskip=10pt,
}

\lstdefinestyle{cbainline}{
  language=cambria,
  basicstyle=\ttfamily\small,
  keywordstyle=\color{cbaKeyword}\bfseries,
  stringstyle=\color{cbaString},
  commentstyle=\color{cbaComment}\itshape,
}

\lstdefinestyle{bash}{
  language=bash,
  basicstyle=\ttfamily\small,
  backgroundcolor=\color{cbaBg},
  frame=single,
  rulecolor=\color{cbaFrame},
  framesep=6pt,
  xleftmargin=8pt,
  xrightmargin=8pt,
  showstringspaces=false,
  aboveskip=10pt,
  belowskip=10pt,
}

\lstset{style=cba}

% Inline code command
\newcommand{\cba}[1]{\lstinline[style=cbainline]{#1}}

% ── Headers / footers ────────────────────────────────────────────────────────
\pagestyle{fancy}
\fancyhf{}
\fancyhead[L]{\small\textit{Cambria Language Reference}}
\fancyhead[R]{\small\thepage}
\renewcommand{\headrulewidth}{0.4pt}

% ── Title tweaks ─────────────────────────────────────────────────────────────
\titleformat{\section}
  {\Large\bfseries\color{cbaKeyword}}{\thesection}{1em}{}
\titleformat{\subsection}
  {\large\bfseries}{\thesubsection}{1em}{}
\titleformat{\subsubsection}
  {\normalsize\bfseries}{\thesubsubsection}{1em}{}

% ══════════════════════════════════════════════════════════════════════════════
\begin{document}

% ── Title page ───────────────────────────────────────────────────────────────
\begin{titlepage}
\centering
\vspace*{4cm}
{\Huge\bfseries Cambria Language Reference\par}
\vspace{1.5cm}
{\Large A Functional Language with\\[4pt]
Parameterized Algebraic Effects and Handlers\par}
\vspace{3cm}
{\large\itshape A complete guide for programming in Cambria\par}
\vfill
\end{titlepage}

% ── Table of contents ────────────────────────────────────────────────────────
\tableofcontents
\newpage

% ══════════════════════════════════════════════════════════════════════════════
\section{Running Programs}
\label{sec:running}

\begin{lstlisting}[style=bash]
cabal build
cabal run cambria -- <file.cba>
\end{lstlisting}

Programs are written in \texttt{.cba} files.  The interpreter parses, desugars,
type-checks, and evaluates the program, printing the final result together with
its inferred type.

Output format: \texttt{<result> : <type>} (e.g.\ \texttt{Pure: 5 : Int!\{\}}).

% ══════════════════════════════════════════════════════════════════════════════
\section{Values and Types}
\label{sec:types}

Cambria distinguishes between \textbf{values} (inert data) and
\textbf{computations} (potentially effectful expressions).  This is a
fine-grained call-by-value discipline.

% ── Primitive types ──────────────────────────────────────────────────────────
\subsection{Primitive Types}

\begin{center}
\begin{tabular}{lll}
\toprule
\textbf{Type} & \textbf{Literal examples} & \textbf{Description} \\
\midrule
\texttt{Unit}   & \cba{()}                    & The unit value \\
\texttt{Int}    & \texttt{0}, \texttt{42}, \texttt{-3}  & Arbitrary-precision integers \\
\texttt{Bool}   & \cba{true}, \cba{false}     & Booleans \\
\texttt{Double} & \texttt{1.0}, \texttt{0.5}, \texttt{-0.25} & Floating-point numbers \\
\texttt{Str}    & \texttt{"hello"}, \texttt{"a\textbackslash nb"} & Strings (with escapes) \\
\texttt{Name} & (produced by \cba{!fresh}) & Opaque unique identifiers \\
\bottomrule
\end{tabular}
\end{center}

String escape sequences: \texttt{\textbackslash n}~(newline),
\texttt{\textbackslash t}~(tab), \texttt{\textbackslash"}~(double quote),
\texttt{\textbackslash\textbackslash}~(backslash).

\subsubsection*{Numeric Literals}

Point notation separates the two numeric types: \texttt{1} is an \texttt{Int}
and \texttt{1.0} is a \texttt{Double}.  A \texttt{Double} literal needs digits
on both sides of the point, so \texttt{1.} and \texttt{.5} are not literals.

A \texttt{-} written directly against the digits belongs to the literal, so
\texttt{-3} and \texttt{(-3, 4)} and \texttt{3 - -2} all contain a negative
literal.  Where an operand comes first there is nothing to negate and the
same text is a subtraction, so \texttt{n-1}, \texttt{n - 1} and \texttt{n -1}
all subtract.

The one thing to know is that a signed literal cannot be an application
argument on its own: \cba{f -3} is the subtraction \cba{f - 3}, and applying
\texttt{f} to \texttt{-3} is written \cba{f (-3)}.

% ── Compound types ───────────────────────────────────────────────────────────
\subsection{Compound Types}

\subsubsection*{Pairs \texttt{A * B}}

\begin{lstlisting}
(1, 2)              -- Int * Int
((1, 2), true)      -- (Int * Int) * Bool
\end{lstlisting}

\subsubsection*{Sum Types (Either) \texttt{A + B}}

\begin{lstlisting}
inl 42              -- creates a left injection
inr true            -- creates a right injection
\end{lstlisting}

\subsubsection*{Lists \texttt{List A}}

\begin{lstlisting}
[]                  -- empty list
1 :: 2 :: 3 :: []   -- list [1, 2, 3]
\end{lstlisting}

\subsubsection*{Maps \texttt{Map K V}}

\begin{lstlisting}
empty                       -- empty map
insert ((key, val), m)      -- insert into map
\end{lstlisting}

% ── Type syntax ──────────────────────────────────────────────────────────────
\subsection{Type Syntax}

Types are written in \cba{effect} declarations and in explicit annotations
(Section~\ref{sec:annotations}), using this syntax:

\begin{center}
\begin{tabular}{ll}
\toprule
\textbf{Syntax} & \textbf{Meaning} \\
\midrule
\texttt{Unit}        & Unit type \\
\texttt{Void}        & Empty type (no values) \\
\texttt{Int}         & Integer type \\
\texttt{Bool}        & Boolean type \\
\texttt{Double}      & Floating-point type \\
\texttt{Str}         & String type \\
\texttt{Name}        & Fresh-name identifier type \\
\texttt{A * B}       & Pair type (product) \\
\texttt{A + B}       & Sum type (either) \\
\texttt{Map K V}     & Map type \\
\texttt{List A}      & List type \\
\texttt{A -> B!\{ops\}} & Function type \\
\texttt{A!\{ops\} => B!\{ops\}} & Handler type \\
\texttt{\$p}         & Abstract type parameter \\
\texttt{(A)}         & Parenthesised type \\
\bottomrule
\end{tabular}
\end{center}

Precedence: \texttt{+} binds looser than \texttt{*}, so
\texttt{Int * Bool + Str} means \texttt{(Int * Bool) + Str}.  Neither operator
associates, so a chain of three or more needs parentheses.

A function type carries the computation type of its result, effects included,
so an operation can take a callback:

\begin{lstlisting}
effect !atomically : (Unit -> Int!{ get : Unit ~> Int }) ~> Int.
\end{lstlisting}

A handler type is written with a fat arrow between two computation types,
the one it consumes and the one it produces (Section~\ref{sec:handlertypes}).

\begin{lstlisting}
(handler { return x -> return x } : Int!{} => Int!{})
\end{lstlisting}

\texttt{Void} has no values and no literal.  It is the sensible result type
for an operation that never returns, such as a jump or a thread halting
(\cba{effect !stop : Unit ~> Void.}). The built-in \cba{absurd} turns a
\texttt{Void} into anything.

% ══════════════════════════════════════════════════════════════════════════════
\section{Values and Computations}
\label{sec:expressions}

Cambria programs are built from two syntactic categories:

\begin{itemize}
  \item \textbf{Values}: inert data that can be passed around
        (integers, booleans, functions, pairs, etc.)
  \item \textbf{Computations}: expressions that may perform effects or produce
        a value
\end{itemize}

A top-level program is a computation.

\subsection{\texttt{return}}

Lifts a value into a computation that does nothing and produces that value:

\begin{lstlisting}
return 42
return (1, 2)
return ()
\end{lstlisting}

\subsection{\texttt{do \ldots\ <- \ldots\ in \ldots}}

Sequences computations, binding the result of the first to a name used in the
second:

\begin{lstlisting}
do x <- return 42 in
do y <- return 10 in
return (x + y)
\end{lstlisting}

The bound name is available in everything after \cba{in}.

\subsection{Semicolon \texttt{;}}

A shorthand for \cba{do _ <- ... in ...} that sequences two computations,
discarding the first result:

\begin{lstlisting}
!print "hello" ; !print "world" ; return ()
\end{lstlisting}

This desugars to:

\begin{lstlisting}
do _ <- !print "hello" in
do _ <- !print "world" in
return ()
\end{lstlisting}

% ══════════════════════════════════════════════════════════════════════════════
\section{Functions}
\label{sec:functions}

\subsection{Anonymous Functions (\texttt{fun})}

\begin{lstlisting}
fun x -> return (x + 1)
\end{lstlisting}

Multi-argument functions are written with multiple variables:

\begin{lstlisting}
fun x y -> return (x + y)
\end{lstlisting}

This desugars to a curried function:
\cba{fun x -> return (fun y -> return (x + y))}.

Function bodies are computations (they can perform effects).

\subsection{Recursive Functions (\texttt{rec})}

\begin{lstlisting}
rec factorial n ->
  if n == 0 then return 1
  else do r <- factorial (n - 1) in
       return (n * r)
\end{lstlisting}

The first identifier after \cba{rec} is the function's own name (for
recursion), followed by argument patterns:

\begin{lstlisting}
rec f x y -> ...    -- f is the recursive name, x and y are arguments
\end{lstlisting}

\subsection{Function Application}

Function application is juxtaposition:

\begin{lstlisting}
f x           -- apply f to x
f (x, y)      -- apply f to the pair (x, y)
\end{lstlisting}

Application is a computation and may trigger effects.  Since functions take
a single argument, multi-argument application is curried:

\begin{lstlisting}
do add <- return (fun x y -> return (x + y)) in
add 1 2
\end{lstlisting}

Here \cba{add 1} returns a closure, and then \texttt{2} is applied to that
closure.

% ══════════════════════════════════════════════════════════════════════════════
\section{Pattern Matching}
\label{sec:patterns}

Patterns appear in \cba{do}-bindings, function arguments, \cba{rec} arguments,
handler clauses, and \cba{case} branches.  A pattern is one of:

\begin{itemize}
  \item A variable name (e.g.\ \texttt{x}): binds the matched value
  \item A wildcard \cba{_}: matches anything, discards the value
  \item A pair pattern \texttt{(\emph{p1}, \emph{p2})}: destructs a pair,
        matching each component with a sub-pattern
\end{itemize}

\subsection{Wildcards}

Use \cba{_} anywhere a pattern is expected to discard a value:

\begin{lstlisting}
do _ <- !print "hello" in return ()
fun _ -> return 42
rec f _ -> return 0
\end{lstlisting}

Wildcards also work inside pair patterns:

\begin{lstlisting}
do (x, _) <- return (1, 2) in return x
fun (_, y) -> return y
\end{lstlisting}

\subsection{Pair Patterns}

Pairs can be destructured wherever patterns are accepted:

\begin{lstlisting}
-- In do-bindings
do (x, y) <- return (1, 2) in return (x + y)

-- In function arguments
fun (x, y) -> return (x + y)

-- Nested patterns
fun ((a, b), c) -> return (a + b + c)

-- In handler operation clauses
set (ref, val) k -> ...
\end{lstlisting}

Pair patterns desugar into calls to the built-in \cba{fst} and \cba{snd}
functions.

% ══════════════════════════════════════════════════════════════════════════════
\section{Control Flow}
\label{sec:control}

\subsection{If-Then-Else}

\begin{lstlisting}
if condition then computation1 else computation2
\end{lstlisting}

The condition is a value of type \texttt{Bool}, and both branches are
computations:

\begin{lstlisting}
if x == 0 then return "zero" else return "nonzero"
\end{lstlisting}

\subsection{Case Expressions (Sum Types)}

Pattern match on sum types (\cba{inl}/\cba{inr}):

\begin{lstlisting}
case expr of {
  inl x -> computation1,
  inr y -> computation2
}
\end{lstlisting}

The two branches can appear in either order:

\begin{lstlisting}
case v of {
  inr y -> return y,
  inl x -> return x
}
\end{lstlisting}

% ══════════════════════════════════════════════════════════════════════════════
\section{Operators}
\label{sec:operators}

\subsection{Arithmetic}

\begin{center}
\begin{tabular}{lll}
\toprule
\textbf{Operator} & \textbf{Types} & \textbf{Result} \\
\midrule
\texttt{+} & \texttt{Int * Int} & \texttt{Int} \\
\texttt{-} & \texttt{Int * Int} & \texttt{Int} \\
\texttt{*} & \texttt{Int * Int} & \texttt{Int} \\
\texttt{/} & \texttt{Int * Int} & \texttt{Double} \\
\bottomrule
\end{tabular}
\end{center}

Note: \texttt{/} always produces a \texttt{Double}, even for integer operands.

\subsection{Comparison}

\begin{center}
\begin{tabular}{lll}
\toprule
\textbf{Operator} & \textbf{Types} & \textbf{Result} \\
\midrule
\texttt{==} & \texttt{a * a}     & \texttt{Bool} \\
\texttt{/=} & \texttt{a * a}     & \texttt{Bool} \\
\texttt{<}  & \texttt{Int * Int} & \texttt{Bool} \\
\texttt{>}  & \texttt{Int * Int} & \texttt{Bool} \\
\texttt{<=} & \texttt{Int * Int} & \texttt{Bool} \\
\texttt{>=} & \texttt{Int * Int} & \texttt{Bool} \\
\bottomrule
\end{tabular}
\end{center}

The \texttt{==} and \texttt{/=} operators are polymorphic. They work on any
type where structural equality is defined (integers, doubles, booleans,
strings, unique values, pairs, lists, maps, sum types).  The four orderings
are on \texttt{Int}, like the other arithmetic.  All six are
non-associative, so \texttt{a < b < c} is a parse error.

\subsection{Boolean Connectives}

\begin{center}
\begin{tabular}{lll}
\toprule
\textbf{Operator} & \textbf{Types} & \textbf{Result} \\
\midrule
\texttt{\&\&} & \texttt{Bool * Bool} & \texttt{Bool} \\
\texttt{||}   & \texttt{Bool * Bool} & \texttt{Bool} \\
\bottomrule
\end{tabular}
\end{center}

\begin{lstlisting}
if x == 0 || 0 < y && y <= 10 then return true else return false
\end{lstlisting}

Both are right-associative, and \cba{||} binds looser than \cba{&&}, so the
condition above reads \cba{(x == 0) || ((0 < y) && (y <= 10))}.  Neither is
short-circuiting: they are ordinary functions on a pair, so both arguments
are evaluated.

\subsection{String Concatenation}

\begin{center}
\begin{tabular}{lll}
\toprule
\textbf{Operator} & \textbf{Types} & \textbf{Result} \\
\midrule
\texttt{++} & \texttt{Str * Str} & \texttt{Str} \\
\bottomrule
\end{tabular}
\end{center}

\subsection{List Cons}

\begin{center}
\begin{tabular}{lll}
\toprule
\textbf{Operator} & \textbf{Types} & \textbf{Result} \\
\midrule
\texttt{::} & \texttt{a * List a} & \texttt{List a} \\
\bottomrule
\end{tabular}
\end{center}

\begin{lstlisting}
1 :: 2 :: 3 :: []    -- builds [1, 2, 3]
\end{lstlisting}

\texttt{::} is right-associative.

\subsection{Operator Precedence (Low to High)}

\begin{enumerate}
  \item \texttt{;} (sequencing): right-associative
  \item \texttt{::} (cons): right-associative
  \item \texttt{||} (disjunction): right-associative
  \item \texttt{\&\&} (conjunction): right-associative
  \item \texttt{==}, \texttt{/=}, \texttt{<}, \texttt{>}, \texttt{<=},
        \texttt{>=} (comparison): non-associative
  \item \texttt{++} (string concat): left-associative
  \item \texttt{+}, \texttt{-} (additive): left-associative
  \item \texttt{*}, \texttt{/} (multiplicative): left-associative
  \item Function application
\end{enumerate}

All infix operators desugar into function applications.  For example,
\texttt{x + y} becomes \texttt{(+) (x, y)}.

The compound computation forms are looser still.  The body of a \cba{do},
the branches of an \cba{if} and a \cba{case}, the body of a
\cba{with ... handle}, and the scope of an \cba{effect} declaration all
extend as far to the right as they can, across \texttt{;} and to the end of
the enclosing computation.  So

\begin{lstlisting}
do x <- return 1 in return x ; return x
\end{lstlisting}

binds \texttt{x} in both computations of the sequence, rather than closing
the \cba{do} at the semicolon.  This is what lets a handler be written as a
prefix, with its body simply following it:

\begin{lstlisting}
with handler { ... } handle

effect !op : Unit ~> Int.
!op () ; !op ()
\end{lstlisting}

Parenthesise a computation to end its scope early.

% ══════════════════════════════════════════════════════════════════════════════
\section{Effects and Operations}
\label{sec:effects}

Effects are the core feature of Cambria.  An \textbf{effect operation} is
invoked with the \texttt{!} prefix:

\begin{lstlisting}
!print "hello"
!flip ()
!get a
!set (a, 42)
\end{lstlisting}

An operation invocation \texttt{!op arg} is a computation that performs the
named effect and produces a result.  The result depends on how the enclosing
handler interprets the operation.

\subsection{Built-in vs User-Defined Operations}

Some operations are built-in and handled by the runtime (see
Section~\ref{sec:builtinio}).  Operations can be:

\begin{enumerate}
  \item \textbf{Declared} with \cba{effect} to enforce a type signature, and
  \item \textbf{Handled} by an enclosing \cba{with ... handle ...} block
\end{enumerate}

If an operation is invoked but no handler catches it, it propagates outward.
If it reaches the top level and is not a built-in, the program reports an
unhandled effect.

% ══════════════════════════════════════════════════════════════════════════════
\section{Handlers}
\label{sec:handlers}

Handlers are the mechanism for defining the meaning of effect operations.  A
handler intercepts operations performed within its body and provides
implementations for them.

\subsection{Handler Syntax}

\begin{lstlisting}
with [$p -> Type] handler {
  return x -> returnBody,
  opName pattern contName -> opBody,
  finally pattern -> finallyBody
} handle

computation
\end{lstlisting}

A handler is a value (of handler type) applied to a computation via
\cba{with ... handle ...}.  The optional \cba{[$p -> Type, ...]} list between
\cba{with} and the handler instantiates the handler's abstract type
parameters (Section~\ref{sec:typeparams}); it is omitted when there are none.

The handled computation runs to the end of the enclosing computation
(Section~\ref{sec:operators}), so a stack of handlers reads as a prefix and
needs no parentheses.  Wrapping the body in \texttt{(\ldots)} is still
allowed, and is how a handler is closed off early.

\subsection{Handler Clauses}

A handler can contain three kinds of clauses, separated by commas. The ordering is not important.

\subsubsection{Return Clause: \texttt{return x -> \ldots}}

Transforms the final value produced by the handled computation:

\begin{lstlisting}
handler {
  return x -> return (x + 1)
}
\end{lstlisting}

If omitted, it defaults to the identity clause: \cba{return x -> return x}.

\subsubsection{Operation Clause: \texttt{opName pattern contName -> \ldots}}

Handles an effect operation.  When \texttt{!opName arg} is invoked inside the
handled computation:

\begin{itemize}
  \item \texttt{pattern} binds the argument \texttt{arg}
  \item \texttt{contName} binds the \textbf{continuation}, which is a function
        representing ``the rest of the computation after the operation''
\end{itemize}

\begin{lstlisting}
handler {
  get a k -> k (lookup (a, s)),
  set x k -> k ()
}
\end{lstlisting}

Calling the continuation \texttt{k} with a value resumes the computation as if
the operation returned that value.  You can:

\begin{itemize}
  \item Call \texttt{k} once (normal resumption)
  \item Call \texttt{k} multiple times (nondeterminism, backtracking)
  \item Not call \texttt{k} at all (abort, short-circuit)
  \item Call \texttt{k} inside another computation (transform the result)
\end{itemize}

\subsubsection{Finally Clause: \texttt{finally f -> \ldots}}

Transforms the entire result of the handler after the return clause has been
applied:

\begin{lstlisting}
handler {
  return x -> return (fun _ -> return x),
  get a k -> return (fun s -> k (lookup (a, s)) s),
  finally s -> s empty
}
\end{lstlisting}

The finally clause is particularly useful for the ``state-passing'' pattern
where the return clause wraps the result in a function, and \cba{finally}
applies that function to an initial state.

\subsection{Handler as a Value}

Handlers are first-class values.  You can bind them to variables:

\begin{lstlisting}
do h <- return (handler {
  return x -> return x,
  decide v k -> max (k true, k false)
}) in
with h handle ...
\end{lstlisting}

\subsection{Nested Handlers}

Handlers can be nested.  The innermost handler for a given operation takes
priority:

\begin{lstlisting}
with outerHandler handle
with innerHandler handle
-- innerHandler's operations are checked first
...
\end{lstlisting}

If an inner handler does not handle an operation, it propagates to the outer
handler.  The inner handler automatically wraps the continuation to maintain
its own handling scope (this is called ``deep handling'').

\subsection{Deep Handling}

When a handler catches an operation and calls the continuation \texttt{k}, the
handler is automatically re-installed around the continuation.  This means
subsequent operations in the same computation are also caught by the same
handler.  This is the standard ``deep handler'' semantics.

% ══════════════════════════════════════════════════════════════════════════════
\section{Type Parameters (\texttt{\$p})}
\label{sec:typeparams}

Type parameters are Cambria's mechanism for \textbf{abstract types}. Types
whose concrete representation is hidden from the code that uses them but known
to the handler that provides them.

\subsection{Declaring Operations with Type Parameters}

\begin{lstlisting}
effect !ref : Int ~> $p.
effect !get : $p ~> Int.
effect !set : $p * Int ~> Unit.
\end{lstlisting}

Here \texttt{\$p} is an abstract type.  Code that uses these operations knows
that \cba{!ref} produces a \texttt{\$p} and \cba{!get} consumes a
\texttt{\$p}, but it cannot inspect or construct \texttt{\$p} values directly.
Attempting to use a \texttt{\$p} value where a concrete type is expected
(e.g.\ \texttt{a + 1} where \texttt{a : \$p}) is a type error.

\subsection{Instantiating Type Parameters in Handlers}

A handler application instantiates a type parameter with a
\texttt{[\$p -> Type]} list placed between \cba{with} and the handler:

\begin{lstlisting}
with [$p -> Name] handler {
  return x -> ...,
  ref x k -> do a <- !fresh () in k a,
  get a k -> ...,
  set x k -> ...
} handle

effect !ref : Int ~> $p.
effect !get : $p ~> Int.
effect !set : $p * Int ~> Unit.
...
\end{lstlisting}

When the handler instantiates \texttt{\$p -> Name}, all occurrences of
\texttt{\$p} in the handled body's effect signatures are replaced with
\texttt{Name} for type-checking purposes.

\subsection{Multiple Type Parameters}

Different parameters can be used for different abstract types:

\begin{lstlisting}
with [$p -> Name, $q -> Int] handler { ... } handle ...
\end{lstlisting}

\subsection{Nested Parameter Instantiation}

A type parameter can be instantiated to a type involving another parameter:

\begin{lstlisting}
with [$p1 -> $p2 + Int] handler {
  ...
} handle ...
\end{lstlisting}

\subsection{Parameter Coverage Checking}

If a handler instantiates a type parameter, it \textbf{must} handle all
operations whose signatures mention that parameter.  Otherwise, the type
checker rejects the program:

\begin{lstlisting}
-- ERROR: handler instantiates $p but does not handle !set
with [$p -> Name] handler {
  return x -> return x,
  get a k -> k 42,
  ref x k -> k ()
} handle

effect !get : $p ~> Int.
effect !set : $p * Int ~> Unit.   -- not handled!
effect !ref : Int ~> $p.
...
\end{lstlisting}

The operations that would be left over are reported by name:
\texttt{Forwarded operations ["set"] mention type parameters bound by this
handler}.  A handler that instantiates \texttt{\$p} is the only place where
\texttt{\$p} has a meaning, so nothing mentioning it may be passed further
out.

\subsection{Parameter Escape Checking}

The same reasoning applies to values, not just to operations.  A
\texttt{\$p}-typed value may not escape the handler that instantiates
\texttt{\$p}.  The type checker rejects a body that leaks one into the surrounding context:

\begin{lstlisting}
-- ERROR: $p escapes through f, which comes from outside the handler
return (fun f ->
  with [$p -> Int] handler {
    return x -> return x,
    get a k -> k 5
  } handle
  effect !get : Unit ~> $p.
  do x <- !get () in
  f x)
\end{lstlisting}

% ══════════════════════════════════════════════════════════════════════════════
\section{Effect Declarations}
\label{sec:effect}

\subsection{\texttt{effect}}

Declares the type signature of an effect operation for use in subsequent code:

\begin{lstlisting}
effect !opName : ArgType ~> RetType.
\end{lstlisting}

The declaration must end with a \texttt{.} (dot) and is followed by the
computation that uses the operation.

Examples:

\begin{lstlisting}
effect !get : $p ~> Int.
effect !set : $p * Int ~> Unit.
effect !ref : Int ~> $p.
effect !ask : Unit ~> Int.
effect !urn : Unit ~> $p.
effect !sample : $p ~> Bool.
\end{lstlisting}

Declarations serve two purposes:

\begin{enumerate}
  \item They tell the type checker the argument and return types of operations.
  \item They declare if the operations use abstract types.
\end{enumerate}

% ══════════════════════════════════════════════════════════════════════════════
\section{Built-in Functions}
\label{sec:builtinfns}

These are available in every program without any declaration.

\subsection{Pair Operations}

\begin{center}
\begin{tabular}{lll}
\toprule
\textbf{Function} & \textbf{Type} & \textbf{Description} \\
\midrule
\texttt{fst} & \texttt{a * b -> a} & First component \\
\texttt{snd} & \texttt{a * b -> b} & Second component \\
\bottomrule
\end{tabular}
\end{center}

\subsection{Arithmetic / Comparison}

\begin{center}
\begin{tabular}{lll}
\toprule
\textbf{Function} & \textbf{Type} & \textbf{Description} \\
\midrule
\texttt{+}   & \texttt{Int * Int -> Int}    & Addition \\
\texttt{-}   & \texttt{Int * Int -> Int}    & Subtraction \\
\texttt{*}   & \texttt{Int * Int -> Int}    & Multiplication \\
\texttt{/}   & \texttt{Int * Int -> Double} & Division \\
\texttt{max} & \texttt{Int * Int -> Int}    & Maximum \\
\texttt{<}   & \texttt{Int * Int -> Bool}   & Less than \\
\texttt{>}   & \texttt{Int * Int -> Bool}   & Greater than \\
\texttt{<=}  & \texttt{Int * Int -> Bool}   & Less than or equal \\
\texttt{>=}  & \texttt{Int * Int -> Bool}   & Greater than or equal \\
\texttt{\&\&} & \texttt{Bool * Bool -> Bool} & Conjunction \\
\texttt{||}  & \texttt{Bool * Bool -> Bool} & Disjunction \\
\texttt{==}  & \texttt{a * a -> Bool}       & Structural equality \\
\texttt{/=}  & \texttt{a * a -> Bool}       & Structural inequality \\
\bottomrule
\end{tabular}
\end{center}

\subsection{String Operations}

\begin{center}
\begin{tabular}{lll}
\toprule
\textbf{Function} & \textbf{Type} & \textbf{Description} \\
\midrule
\texttt{++}   & \texttt{Str * Str -> Str} & Concatenation \\
\texttt{hash} & \texttt{Name -> Str}      & Convert a name to a string \\
\bottomrule
\end{tabular}
\end{center}

\subsection{List Operations}

\begin{center}
\begin{tabular}{lll}
\toprule
\textbf{Function} & \textbf{Type} & \textbf{Description} \\
\midrule
\texttt{::}     & \texttt{a * List a -> List a}        & Prepend element \\
\texttt{uncons} & \texttt{List a -> Unit + a * List a} & List destructor \\
\texttt{null}   & \texttt{List a -> Bool}               & Is empty? \\
\texttt{[]}     & \texttt{List a}                       & Empty list constant \\
\bottomrule
\end{tabular}
\end{center}

\subsection{Map Operations}

\begin{center}
\begin{tabular}{lll}
\toprule
\textbf{Function} & \textbf{Type} & \textbf{Description} \\
\midrule
\texttt{empty}  & \texttt{Map k v}                       & Empty map constant \\
\texttt{insert} & \texttt{(k * v) * Map k v -> Map k v} & Insert or update entry \\
\texttt{remove} & \texttt{k * Map k v -> Map k v}        & Remove entry by key \\
\texttt{lookup} & \texttt{k * Map k v -> v}              & Lookup by key (partial) \\
\texttt{member} & \texttt{k * Map k v -> Bool}           & Key exists? \\
\bottomrule
\end{tabular}
\end{center}

Note: \texttt{lookup} will error at runtime if the key is not found.
Check with \texttt{member} first if unsure.

\subsection{The Empty Type}

\begin{center}
\begin{tabular}{lll}
\toprule
\textbf{Function} & \textbf{Type} & \textbf{Description} \\
\midrule
\texttt{absurd} & \texttt{Void -> a} & Anything follows from \texttt{Void} \\
\bottomrule
\end{tabular}
\end{center}

An operation declared to return \texttt{Void} never resumes, so its result
can be turned into whatever the surrounding code needs:

\begin{lstlisting}
effect !stop : Unit ~> Void.
do v <- !stop () in absurd v
\end{lstlisting}

% ══════════════════════════════════════════════════════════════════════════════
\section{Built-in IO Effects}
\label{sec:builtinio}

These operations are handled at the top level by the runtime and perform actual
I/O:

\begin{center}
\begin{tabular}{lll}
\toprule
\textbf{Operation} & \textbf{Signature} & \textbf{Description} \\
\midrule
\texttt{!fresh ()}    & \texttt{Unit $\leadsto$ Name}    & Fresh unique identifier \\
\texttt{!print s}      & \texttt{Str $\leadsto$ Unit}       & Print string to stdout \\
\texttt{!read ()}      & \texttt{Unit $\leadsto$ Str}       & Read line from stdin \\
\texttt{!flip ()}      & \texttt{Unit $\leadsto$ Bool}      & Random boolean (fair coin) \\
\texttt{!bernoulli p}  & \texttt{Double $\leadsto$ Bool}    & Random boolean with probability \texttt{p} \\
\texttt{!uniform ()}   & \texttt{Unit $\leadsto$ Double}    & Random double in $[0, 1)$ \\
\bottomrule
\end{tabular}
\end{center}

These do not need to be declared; they are always available.  They are
caught at the outermost level after all user-defined handlers have had a
chance to intercept them.

% ══════════════════════════════════════════════════════════════════════════════
\section{The Type System}
\label{sec:typesystem}

Cambria uses an extended Hindley-Milner type inference system with effect types.
The only annotation that must be supplied is the arity of any operation whose signature mentions an abstract parameter type (\texttt{\$p}),
which are introduced with an \cba{effect} declaration.
Explicit annotations are nonetheless available where desired.

\subsection{Computation Types}

Every computation has a type of the form:

\begin{center}
\texttt{ValueType!\{effects\}}
\end{center}

For example:

\begin{itemize}
  \item \texttt{Int!\{\}}: a pure computation producing an \texttt{Int}
  \item \texttt{Bool!\{ flip : Unit \(\leadsto\) Bool \}}: a computation
        that may invoke \texttt{!flip} and produces a \texttt{Bool}
  \item \texttt{Unit!\{ print : Str \(\leadsto\) Unit, fresh : Unit
        \(\leadsto\) Name \}}: may print and generate uniques
\end{itemize}

\subsection{Effect Rows}

Effects are tracked as a set of operation signatures.  When operations are
handled, they are removed from the effect set.  A fully handled program has an
empty effect set \texttt{\{\}} (plus any built-in IO effects that are handled
by the runtime).

Effect rows can be either \textbf{closed} (a fixed set of operations) or
\textbf{open} (a set of known operations plus an effect variable representing
unknown additional effects).  Open effect rows allow polymorphism over effects.

The two print differently.  A closed row is just its operations between
braces, and an open row adds a bar and the effect variable:

\begin{center}
\begin{tabular}{ll}
\toprule
\textbf{Printed} & \textbf{Meaning} \\
\midrule
\texttt{\{\}}                        & Closed and empty: no effects at all \\
\texttt{\{ print : Str \(\leadsto\) Unit \}} & Closed: exactly this operation \\
\texttt{\{\_e1\}}                    & Open and empty: any effects \\
\texttt{\{ print : Str \(\leadsto\) Unit | \_e1 \}} & Open: this operation and possibly more \\
\bottomrule
\end{tabular}
\end{center}

\subsection{Polymorphism}

Functions defined with \cba{do ... <- return (fun ...) in ...} or
\cba{do ... <- return (rec ...) in ...} are generalised so their types are
polymorphic. The generalisation happens at the \cba{do}-binding site.

\begin{lstlisting}
do id <- return (fun x -> return x) in
-- id has type: forall a. a -> a!{e}
do _ <- id 42 in         -- instantiated at Int
do _ <- id true in        -- instantiated at Bool
return ()
\end{lstlisting}

Recursive functions are also polymorphic:

\begin{lstlisting}
do map <- return (rec map f xs ->
  case uncons xs of {
    inl _        -> return [],
    inr (x, xs') -> f x :: map f xs'
  }
) in
-- map : forall a b. (a -> b!S) -> List a -> List b!S
\end{lstlisting}

\subsection{Type Parameter Abstraction}

Type parameters (\texttt{\$p}) are existentially quantified from the body's
perspective.  The body knows operations involve \texttt{\$p} but cannot see its
concrete type.  Only the handler knows the instantiation.

This enables abstraction: you can write polymorphic functions that work over
\texttt{\$p} without knowing what it is:

\begin{lstlisting}
-- This swap function works at both Int and $p
do swap <- return (fun (x, y) -> return (y, x)) in
swap (1, 2)              -- at Int
do a <- !ref 100 in
do b <- !ref 200 in
swap (a, b)              -- at $p (Name, but swap doesn't know)
\end{lstlisting}

\subsection{Handler Types}
\label{sec:handlertypes}

Handlers have types of the form:

\begin{center}
\texttt{CompType\textsubscript{in} => CompType\textsubscript{out}}
\end{center}

The input computation type includes the effects the handler catches.  The
output computation type is the type of the result after handling.  Both rows
are usually open, so the collecting handler of Section~\ref{sec:examples} has
a type of the shape

\begin{center}
\texttt{(a!\{ print : Str \(\leadsto\) Unit | e \} => (a * Str)!\{e'\})}
\end{center}

A handler type records no parameter instantiation.  The
\texttt{[\$p -> Type]} list belongs to the \cba{with ... handle}
application rather than to the handler value, so it is checked where the
handler is applied, not where it is written.

Written out in an annotation, the fat arrow separates the two computation
types.  The handler below catches \cba{!flip} and so discharges it, turning
an \texttt{Int} computation that may flip into a pure one:

\begin{lstlisting}
do h <- return ((handler {
  return x -> return x,
  flip _ k -> k true
}) : Int!{flip : Unit ~> Bool} => Int!{}) in
with h handle
do b <- !flip () in
if b then return 1 else return 2
\end{lstlisting}

% ══════════════════════════════════════════════════════════════════════════════
\section{Explicit Annotations}
\label{sec:annotations}

Cambria infers types without help, so annotations are never required except
for the arity of an operation whose signature mentions a parameter type
(Section~\ref{sec:effect}).  They are available as a way to pin
a type or document intent.

An annotation is a parenthesised expression followed by \texttt{:} and a
type.

\subsection{Value Annotations}

\begin{lstlisting}
return (5 : Int)
do x <- return 7 in return (x : Int)
return ((inl 5) : Int + Bool)
\end{lstlisting}

The annotation sits on the expression, not on the surrounding computation,
which is why the second line annotates \texttt{x} inside the \cba{return}
rather than the \cba{do}-binding.  Annotating a polymorphic expression pins
its instance. Without the annotation, \cba{inl 5} would leave the right-hand
summand undetermined.

A value annotation may be a function type, in which case the result is a
computation type and carries a row:

\begin{lstlisting}
do f <- return ((fun x -> return x) : Int -> Int!{}) in f 4
\end{lstlisting}

Annotating that same function \texttt{Int -> Bool!\{\}} is a type error,
because the body returns its argument.

\subsection{Computation Annotations}

A computation type is a value type, \texttt{!}, and a braced row:

\begin{lstlisting}
(return 3 : Int!{})
(return (1, true) : Int * Bool!{})
(!print "hello" : Unit!{print : Str ~> Unit})
\end{lstlisting}

The row is checked as well as the value type.  A row that omits an effect the
body actually performs is rejected, so \texttt{(!print "boom" : Unit!\{\})}
fails with an effect mismatch, as do two nested annotations that disagree.

\subsection{Open Rows in Annotations}

A row written out in full is \textbf{closed}.  The computation may only perform
the effects it lists.  Ending the row with \texttt{...} makes it
\textbf{open}, so other operations are permitted:

\begin{lstlisting}
(do _ <- !print "hi" in !fresh () : Name!{print : Str ~> Unit, ...})
\end{lstlisting}

The body also performs \cba{!fresh}, which the closed form would reject; the
ellipsis lets the row carry it, and the inferred type comes out with both
operations.

The ellipsis stands for the row variable of Section~\ref{sec:typesystem}, so
an open annotation constrains the row from below rather than fixing it.

% ══════════════════════════════════════════════════════════════════════════════
\section{Desugaring}
\label{sec:desugaring}

The parser produces a ``sugared'' AST that is desugared into a simpler core AST
before type-checking and evaluation.  Understanding desugaring helps explain
how syntactic sugar works.

\subsection{Multi-Argument Functions}

\begin{lstlisting}
fun x y z -> body
\end{lstlisting}

desugars to:

\begin{lstlisting}
fun x -> return (fun y -> return (fun z -> body))
\end{lstlisting}

Each additional argument becomes a nested function that returns the next
closure.

\subsection{Multi-Argument Rec}

\begin{lstlisting}
rec f x y -> body
\end{lstlisting}

desugars to:

\begin{lstlisting}
rec f x -> return (fun y -> body)
\end{lstlisting}

\subsection{Pair Patterns}

\begin{lstlisting}
do (x, y) <- comp in body
\end{lstlisting}

desugars to:

\begin{lstlisting}
do _tmp <- comp in
do x <- fst _tmp in
do y <- snd _tmp in
body
\end{lstlisting}

Nested pair patterns work recursively:

\begin{lstlisting}
fun ((a, b), c) -> body
\end{lstlisting}

desugars to a function on a fresh variable, with projections extracted via
\cba{fst} and \cba{snd}.

\subsection{Infix Operators}

All infix operators desugar to function applications on pairs:

\begin{center}
\begin{tabular}{lcl}
\texttt{x + y}  & $\longrightarrow$ & \texttt{(+) (x, y)} \\
\texttt{x == y} & $\longrightarrow$ & \texttt{(==) (x, y)} \\
\texttt{x ++ y} & $\longrightarrow$ & \texttt{(++) (x, y)} \\
\texttt{x :: y} & $\longrightarrow$ & \texttt{(::) (x, y)} \\
\end{tabular}
\end{center}

\subsection{Semicolons}

\begin{lstlisting}
comp1 ; comp2
\end{lstlisting}

desugars to:

\begin{lstlisting}
do _ <- comp1 in comp2
\end{lstlisting}

\subsection{Inline Computations in Value Position}

When a computation appears where a value is expected (e.g.\ as a function
argument), it is lifted into a \cba{do}-binding:

\begin{lstlisting}
f (!get a)
\end{lstlisting}

desugars to:

\begin{lstlisting}
do _tmp <- !get a in f _tmp
\end{lstlisting}

This is called \textbf{ANF conversion} (A-normal form) and happens
automatically.

\subsection{Default Handler Clauses}

If no \cba{return} clause is given, the default is
\cba{return x -> return x}.

% ══════════════════════════════════════════════════════════════════════════════
\section{Evaluation Model}
\label{sec:eval}

Cambria uses \textbf{call-by-value} evaluation with \textbf{deep effect
handlers}.

\subsection{Values vs Computations at Runtime}

\begin{itemize}
  \item \textbf{Values} are fully evaluated data: integers, closures, pairs,
        etc.
  \item \textbf{Computations} either produce a \textbf{pure result} or an
        \textbf{impure result} (an unhandled operation).
\end{itemize}

The evaluator returns one of:

\begin{itemize}
  \item \texttt{Pure v}: the computation finished with value \texttt{v}
  \item \texttt{Impure op arg}: the computation performed
        operation \texttt{!op arg} and is suspended, waiting for a result via
        the continuation
\end{itemize}

\subsection{Function Evaluation}

\begin{itemize}
  \item \cba{fun x -> body} creates a \textbf{closure} capturing the current
        environment.
  \item \cba{rec f x -> body} creates a closure with a self-reference for
        recursion (via a knot-tying \texttt{let rec}).
  \item Application substitutes the argument into the closure's environment and
        evaluates the body.
\end{itemize}

\subsection{Handler Evaluation}

When \cba{with h handle c} is evaluated:

\begin{enumerate}
  \item Evaluate \texttt{c} in the current environment.
  \item If \texttt{c} produces \texttt{Pure v}:
    \begin{itemize}
      \item Apply the handler's return clause: evaluate the return body with
            \texttt{v} bound to the return pattern variable.
    \end{itemize}
  \item If \texttt{c} produces \texttt{Impure op arg k}:
    \begin{itemize}
      \item If the handler has a clause for \texttt{op}: evaluate the clause
            body with \texttt{arg} bound to the parameter pattern and
            \texttt{k} bound to a \textbf{deep-handled continuation} (a new
            closure that re-wraps the continuation with the same handler).
      \item If the handler does not handle \texttt{op}: propagate the impure
            result upward, wrapping the continuation so the handler remains
            installed.
    \end{itemize}
\end{enumerate}

The \textbf{deep handling} of continuations means that when an operation clause
calls \texttt{k value}, the handler is automatically re-installed around the
remaining computation.

\subsection{Finally Clause Evaluation}

The \cba{finally} clause behaves as a \cba{do}-binding around the whole handler application, so

\begin{lstlisting}
with handler { ..., finally f -> body } handle c
\end{lstlisting}

evaluates like

\begin{lstlisting}
do f <- (with handler { ... } handle c) in body
\end{lstlisting}

This is why the finally clause sees the fully-processed result of the handler.

\subsection{State-Passing Pattern}

The most common handler pattern threads state through continuations:

\begin{lstlisting}
with handler {
  return x  -> return (fun _ -> return x),
  get a k   -> return (fun s -> k (lookup (a, s)) s),
  set x k   -> return (fun s -> k () (insert (x, s))),
  ref x k   -> do a <- !fresh () in
                return (fun s ->
                  k a (insert (((a, x), s)))),
  finally s -> s empty
} handle (...)
\end{lstlisting}

How it works:

\begin{enumerate}
  \item \textbf{return} wraps the final value in a function that ignores state.
  \item \textbf{get} returns a function that reads from state and passes it to
        the continuation.
  \item \textbf{set} returns a function that updates state and passes
        \texttt{()} to the continuation.
  \item \textbf{ref} allocates a fresh unique ID and returns a function that
        extends the state.
  \item \textbf{finally} applies the resulting state-function to an initial
        empty map.
\end{enumerate}

The result of the handler is a function
$\mathit{State} \to \mathit{Result}$.  The \cba{finally} clause seeds it with
\cba{empty}.

% ══════════════════════════════════════════════════════════════════════════════
\section{Complete Examples}
\label{sec:examples}

\subsection{Hello World}

\begin{lstlisting}
!print "Hello, world!"
\end{lstlisting}

\subsection{Local Mutable State}

\begin{lstlisting}
with [$name -> Int] handler {
  return x   -> return (fun _ -> return x),
  fresh _ k  -> return (fun n -> k n (n+1)),
  eq (a,b) k -> k (a == b),
  finally f  -> f 0
} handle

effect !fresh : Unit ~> $name.
with [$loc -> $name] handler {
  return x    -> return (fun _ -> return x),
  get a k     -> return (fun s -> k (s a) s),
  set (a,x) k -> return (fun s ->
                   k () (fun b -> if !eq (a, b) then return x else s b)),
  ref x k     -> do a <- !fresh () in return (fun s ->
                   k a (fun b -> if !eq (a, b) then return x else s b)),
  finally s   -> s (fun _ -> return 0)
} handle

effect !ref : Int ~> $loc.
effect !get : $loc ~> Int.
effect !set : $loc * Int ~> Unit.
do a <- !ref 2 in
do b <- !ref 3 in
do _ <- !set (a, !get b) in
!get a
\end{lstlisting}

The state is encoded as a function from locations to values, and a separate
outer handler supplies the fresh names used as locations.  Each handler is a
prefix of the computation it handles, so no parentheses are needed.  The
program allocates two references (holding 2 and 3), copies \texttt{b}'s value
into \texttt{a}, and reads \texttt{a} back.  Result: \texttt{Pure: 3}.  This
is \texttt{examples/local\_state.cba}.

\subsection{Nondeterministic Maximisation}

\begin{lstlisting}
with handler {
  decide v k -> max (k true, k false)
} handle

do x1 <- if !decide () then return 15 else return 30 in
do x2 <- if !decide () then return 5  else return 10 in
return (x1 - x2)
\end{lstlisting}

The handler explores all branches and returns the maximum.  The \texttt{decide}
operation calls the continuation twice: once with \cba{true}, once with
\cba{false}.  Result: \texttt{Pure: 25} ($30 - 5$).

\subsection{Collecting Output}

\begin{lstlisting}
with handler {
  return x -> return (x, ""),
  print s k -> do (v, acc) <- k () in
    return (v, s ++ acc)
} handle

!print "A"; !print "B"; !print "C"
\end{lstlisting}

Instead of printing to stdout, this handler collects all printed strings into
an accumulator.  Result: \texttt{Pure: ((), "ABC")}.

\subsection{Polymorphic Functions with Abstract Types}

\begin{lstlisting}
with [$name -> Int] handler {
  return x   -> return (fun _ -> return x),
  fresh _ k  -> return (fun n -> k n (n+1)),
  eq (a,b) k -> k (a == b),
  finally f  -> f 0
} handle

effect !fresh : Unit ~> $name.
with [$loc -> $name] handler {
  return x    -> return (fun _ -> return x),
  get a k     -> return (fun s -> k (s a) s),
  set (a,x) k -> return (fun s ->
                   k () (fun b -> if !eq (a, b) then return x else s b)),
  ref x k     -> do a <- !fresh () in return (fun s ->
                   k a (fun b -> if !eq (a, b) then return x else s b)),
  finally s   -> s (fun _ -> return 0)
} handle

effect !ref : Int ~> $loc.
effect !get : $loc ~> Int.
effect !set : $loc * Int ~> Unit.

-- Polymorphic map: forall a b. (a -> b!S) -> List a -> List b!S
do map <- return (rec map f xs ->
  case uncons xs of {
    inl _        -> return [],
    inr (x, xs') -> f x :: map f xs'
  })
in
do vals <- return (10 :: 20 :: 30 :: []) in

-- (1) map at (Int -> $loc): allocate a ref for each number
do refs <- map (fun n -> !ref n) vals in

-- (2) map at ($loc -> Bool): read each ref and test it against 20
do checks <- map (fun n -> !get n == 20) refs in

return checks
\end{lstlisting}

The same polymorphic \texttt{map} is used at two different type
instantiations: \texttt{Int -> \$loc}, allocating a reference per element,
and \texttt{\$loc -> Bool}, reading each reference and comparing it.
\texttt{map} never inspects what \texttt{\$loc} is.  Result:
\texttt{Pure: [False, True, False]}.  This is
\texttt{examples/poly\_map.cba}.

\subsection{Two Implementations of a Probabilistic Urn}

A key strength of parametrized handlers is that the \emph{same} abstract
interface can be given completely different implementations by swapping the
handler.  The following two programs both implement P\'{o}lya's urn with the
same body code and the same effect interface:

\begin{lstlisting}
effect !newurn : Unit ~> $urn.
effect !sample : $urn ~> Bool.
\end{lstlisting}

\texttt{!newurn} creates a new urn and \texttt{!sample} draws a boolean sample
from it.  The abstract type \texttt{\$urn} hides what an urn actually is.  The
shared body draws two samples and returns them as a pair:

\begin{lstlisting}
do urn <- !newurn () in
do b1 <- !sample urn in
do b2 <- !sample urn in
return (b1, b2)
\end{lstlisting}

\subsubsection*{Implementation 1: Real-number parameter (\texttt{\$urn -> Double})}

The simplest implementation instantiates \texttt{\$urn} to \texttt{Double}.  An
urn is just a bias drawn once from the uniform distribution, and each sample is
an independent Bernoulli trial with that bias:

\begin{lstlisting}
with [$urn -> Double] handler {
  newurn _ k -> k (!uniform ()),
  sample p k -> k (!bernoulli p)
} handle

effect !newurn : Unit ~> $urn.
effect !sample : $urn ~> Bool.
do urn <- !newurn () in
do b1 <- !sample urn in
do b2 <- !sample urn in
return (b1, b2)
\end{lstlisting}

The handler is flat: \texttt{newurn} draws a \texttt{Double} and \texttt{sample}
passes it directly to \texttt{!bernoulli}.  No state is needed.  Result: a pair
of booleans, e.g.\ \texttt{Pure: (False, True)}.  This is
\texttt{examples/prob\_polya\_bias.cba}.

\subsubsection*{Implementation 2: P\'{o}lya urn via local state (\texttt{\$urn -> \$loc})}

The P\'{o}lya urn implementation instantiates \texttt{\$urn} to a location
\texttt{\$loc} of the local-state effect, storing a pair of counts (red,
blue).  Each sample draws a colour with probability proportional to the counts
and then adds a ball of that colour, so the urn is biased towards the colour
already drawn (``rich get richer'').  It reuses the function-encoded local
state from above, with its fresh-name handler outside that in turn:

\begin{lstlisting}
with [$name -> Int] handler {
  return x   -> return (fun _ -> return x),
  fresh _ k  -> return (fun n -> k n (n+1)),
  eq (a,b) k -> k (a == b),
  finally f  -> f 0
} handle

effect !fresh : Unit ~> $name.
with [$loc -> $name] handler {
  return x    -> return (fun _ -> return x),
  get a k     -> return (fun s -> k (s a) s),
  set (a,x) k -> return (fun s ->
                   k () (fun b -> if !eq (a, b) then return x else s b)),
  ref x k     -> do a <- !fresh () in return (fun s ->
                   k a (fun b -> if !eq (a, b) then return x else s b)),
  finally s   -> s (fun _ -> return (0, 0))
} handle

effect !ref : Int * Int ~> $loc.
effect !get : $loc ~> Int * Int.
effect !set : $loc * (Int * Int) ~> Unit.
with [$urn -> $loc] handler {
  newurn _ k -> k (!ref (1, 1)),
  sample urn k ->
    do (red, blue) <- !get urn in
    do isRed <- !bernoulli (red / (red + blue)) in
    do _ <- if isRed then !set (urn, (red + 1, blue))
                     else !set (urn, (red, blue + 1)) in
    k isRed
} handle

effect !newurn : Unit ~> $urn.
effect !sample : $urn ~> Bool.
do urn <- !newurn () in
do b1 <- !sample urn in
do b2 <- !sample urn in
return (b1, b2)
\end{lstlisting}

Several things to note:

\begin{itemize}
  \item The \textbf{innermost handler} implements the urn on top of local
        state: \texttt{newurn} allocates a reference holding the initial counts
        \texttt{(1, 1)}, and \texttt{sample} reads the counts, draws a biased
        coin, updates the counts, and returns the result.
  \item It instantiates \texttt{\$urn -> \$loc}: the urn parameter becomes a
        location of the \emph{enclosing} local-state handler, itself an
        abstract parameter.
  \item The \textbf{middle handler} provides the function-encoded local state
        (\texttt{\$loc -> \$name}), and the \textbf{outermost handler} supplies
        the fresh location names (\texttt{\$name -> Int}).
  \item The \textbf{body code is identical} to Implementation~1.  Only the
        handler stack changes.
\end{itemize}

This is \texttt{examples/prob\_polya\_state.cba}.

Both handlers induce the same distribution on the returned pair (P\'{o}lya's
identity), so by parametricity the body cannot tell them apart.

% ══════════════════════════════════════════════════════════════════════════════
\section{Grammar Summary}
\label{sec:grammar}

\begin{lstlisting}[language={},basicstyle=\ttfamily\small,backgroundcolor=\color{cbaBg},frame=single,rulecolor=\color{cbaFrame},framesep=6pt,xleftmargin=8pt,xrightmargin=8pt,keywordstyle={},commentstyle={}]
program := computation

-- Values
value   := atom
         | inl atom | inr atom
         | fun patterns -> computation
         | rec name patterns -> computation
         | handler { clauses }

atom    := () | [] | integer | double
         | boolean | string
         | (expr, expr) | name

integer := digit+
double  := digit+ . digit+

-- A '-' glued to the digits is part of the
-- literal wherever an expression may start,
-- and a subtraction after an operand.  An
-- application argument takes parentheses:
-- f (-3), since f -3 is the subtraction f - 3.

-- Explicit annotations
expr    := ( expr : type )
comp    := ( comp : compType )

-- Computations
comp    := comp ; comp
         | return expr
         | do pattern <- comp in comp
         | if expr then comp else comp
         | case expr of { inl p -> comp,
                          inr p -> comp }
         | with paramInsts expr handle comp
         | effect !op : type ~> type . comp
         | expr expr
         | !op expr
         | expr BINOP expr

BINOP   := :: | || | && | == | /= | < | >
         | <= | >= | ++ | + | - | * | /

-- The trailing comp of do, if, case, with and effect
-- extends as far to the right as it can, across ';'.
-- Parenthesise a comp to end its scope early.

-- Patterns
pattern := name | _ | (pattern, pattern)

-- Handler clauses
clause  := return pattern -> comp
         | opName pattern contName -> comp
         | finally pattern -> comp

-- Parameter instantiations (optional, between with and the handler)
paramInsts := [ paramList ]
paramList  := $param -> type
            | $param -> type , paramList

-- Types (in declarations and annotations)
type     := typeSum
          | typeSum -> compType
          | compType => compType
typeSum  := typeProd
          | typeProd + typeProd
typeProd := typeAtom
          | typeAtom * typeAtom
typeAtom := Unit | Int | Bool | Double
          | Str | Name | Void
          | Map typeAtom typeAtom
          | List typeAtom
          | $param | (type)

compType := type ! { effList }
effList  := empty | ...
          | opSig | opSig , effList
opSig    := opName : type ~> type

-- A trailing ... makes the row open.
\end{lstlisting}

% ══════════════════════════════════════════════════════════════════════════════
\section{Comments}
\label{sec:comments}

Line comments start with \texttt{-{}-}:

\begin{lstlisting}
-- This is a comment
return 42  -- inline comment
\end{lstlisting}

There are no block comments.

\end{document}
